\subsection{Biological Characteristics of \gls{lusc} Subtypes}

%TODO: i think there's sth missing here; doesn't feel substantial, idk what im missing, but i feel like im missing a significant portion here.. like pathways, and things like that, maybe i should also have pathway diagrams key molecule, pathways, axes, clinical severity and things like that.. maybe later

The four established \gls{lusc} molecular subtypes reflect distinct biological axes of squamous tumour identity and understanding these programs is crucial for assessing whether machine learning features recover biologically coherent signals.

\textbf{Classical subtype.}
The classical subtype is enriched for xenobiotic metabolism and oxidative stress response, driven primarily by the NFE2L2/KEAP1 pathway. \textit{NFE2L2} encodes NRF2, a transcription factor that activates antioxidant response element (ARE)-driven genes including those for glutathione synthesis and drug efflux; \textit{KEAP1} normally targets NRF2 for proteasomal degradation, and loss-of-function alterations in \textit{KEAP1} or gain-of-function mutations in \textit{NFE2L2} are recurrently observed in classical \gls{lusc} \parencite{Solis2012, Hammerman2012}. SQSTM1/p62 bridges NRF2 activation and autophagy signalling and is a characteristic protein-level marker of this subtype \parencite{Stewart2019}. TGM1 (transglutaminase~1) encodes an enzyme essential for cornified envelope crosslinking during terminal squamous differentiation and is a well-established transcriptional marker of the classical subtype \parencite{Wilkerson2010}. KRT75 (keratin~75) is a type~II cytoskeletal keratin expressed in differentiated squamous epithelium, complementing the cornification program encoded by TGM1 \parencite{Wilkerson2010, Hammerman2012}.


\textbf{Basal subtype.}
The basal subtype reflects undifferentiated squamous/basal epithelial identity and cell-adhesion programs. TP63 is the master transcriptional regulator of squamous lineage commitment and is characteristically expressed in basal-cell tumours \parencite{Wilkerson2010, Hammerman2012}. Basal keratins including KRT5 and KRT14 mark a less terminally differentiated but committed basal state. NF1 loss-of-function alterations are enriched in this subtype and may activate RAS-MAPK signalling pathways \parencite{Hammerman2012}. The basal subtype presents the greatest classification difficulty in multi-class models because its biological programs overlap partially with both classical (squamous lineage) and primitive (proliferative) subtypes.
 
\textbf{Secretory subtype.}
The secretory subtype is characterised by immune, inflammatory, and secretory epithelial transcriptional programs. Hallmark markers include the innate immune alarmins S100A8 and S100A9, the neutrophil-recruiting chemokine CXCL1, and the cytokine IL-1A \parencite{Wilkerson2010}. The secretory subtype consistently shows the highest immune microenvironment infiltration scores among the four subtypes across multiple studies \parencite{Yang2022immune, Wilkerson2010}, making it of particular interest for understanding immunotherapy response patterns. Claudin-family tight-junction proteins and goblet-cell markers are also expressed in this subtype, reflecting a secretory epithelial cell state.
 
\textbf{Primitive subtype.}
The primitive subtype is associated with high proliferative activity, a less differentiated transcriptional state, and the worst prognosis among the four subtypes \parencite{Wilkerson2010, Hammerman2012}. Recurrent \textit{RB1} and \textit{PTEN} alterations are characteristic, consistent with deregulated cell-cycle control and PI3K/AKT pathway
activation \parencite{Hammerman2012}. EZH2, a Polycomb group histone methyltransferase that catalyses trimethylation of H3K27 (H3K27me3) to silence developmental and differentiation genes, is active in the primitive subtype \parencite{Varambally2008}. hsa-miR-101-3p has been reported to repress EZH2 expression post-transcriptionally in cancer contexts \parencite{Varambally2008, Montani2024}, providing a mechanistic link between the miRNA expression profiles and the epigenetic programs characterising this subtype. hsa-miR-205-5p is a well-characterised squamous-lineage miRNA marker, and its abundance has been shown to distinguish squamous from non-squamous \gls{nsclc} histology with high specificity \parencite{Lebanony2009}.
 
These biological associations define concrete expectations for interpreting machine learning features recovered from each modality. A biologically coherent model should recover oxidative stress and differentiation markers (TGM1, KRT75, NRF2/KEAP1 axis) in the classical subtype space, immune and chemokine signatures (S100A8, CXCL1) in the secretory space, and epigenetic/proliferative markers (EZH2, RB1) in the primitive space. 
